\subsection{ON--OFF 模型的突发流量描述}

为刻画具有突发性的到达过程，本文采用典型的 ON--OFF
二态流量模型。系统在每一个时隙处于 ON 状态或 OFF 状态，
并在不同状态下具有不同的到达率。设

\begin{itemize}
	\item ON 状态时到达率为 $\lambda_1 = 1$，即每个时隙均到达一个数据包；
	\item OFF 状态时到达率为 $\lambda_0 \leq \lambda_1$，其中 $\lambda_0 \in [0,1)$，表示进入常规阶段时的稀疏到达强度；
	\item ON 段持续时间服从参数为 $\sigma_1$ 的几何分布，其均值为 $1/\sigma_1$；
	\item OFF 段持续时间服从参数为 $\sigma_0$ 的几何分布，其均值为 $1/\sigma_0$。
\end{itemize}

因此，总的到达过程可表示为
\[
A_t =
\begin{cases}
	1, & \text{若处于 ON 状态；}\\[3pt]
	\text{Bernoulli}(\lambda_0), & \text{若处于 OFF 状态。}
\end{cases}
\]

该模型能够灵活地控制突发程度：
较大的 ON 段平均长度 ($1/\sigma_1$ 较大) 会导致更长时间的高到达率，从而增强突发性；
较小的 OFF 段平均长度 ($1/\sigma_0$ 较小) 使系统更频繁处于 ON 状态，同样会提升突发性。

为了刻画突发性强弱，可采用经典的离散度指数（Index of Dispersion, ID），其定义为
\[
\mathrm{ID} = \frac{\operatorname{Var}(A_t)}{\mathbb{E}[A_t]}.
\]
当 $\mathrm{ID} > 1$ 时，说明到达过程具有突发性；
当 $\mathrm{ID} = 1$ 时，与泊松过程一致；
当 $\mathrm{ID} < 1$ 时，表示流量较为规则。
该指标能够有效反映 ON--OFF 模型中由状态切换导致的到达率波动，是刻画突发性的自然度量。


\subsection{ON--OFF 模型的突发流量及 ID 闭式推导}

考虑离散时隙的 ON--OFF 模型。设 $S_t \in \{0,1\}$ 表示第 $t$ 个时隙的状态，其中 $S_t=1$ 表示 ON 状态（到达率为 $\lambda_1$），$S_t=0$ 表示 OFF 状态（到达率为 $\lambda_0$）。ON 段和 OFF 段的持续时间分别服从几何分布，其参数为 $\sigma_1$ 和 $\sigma_0$，即每个时隙从 ON 转为 OFF 的概率为 $\sigma_1$，从 OFF 转为 ON 的概率为 $\sigma_0$。

\paragraph{状态的平稳分布}
两态马尔可夫链的平稳概率为：
\[
\pi_{\mathrm{on}} = \frac{\sigma_0}{\sigma_0 + \sigma_1}, \qquad
\pi_{\mathrm{off}} = 1 - \pi_{\mathrm{on}} = \frac{\sigma_1}{\sigma_0 + \sigma_1}.
\]

\paragraph{到达过程的均值}
令 $X_t$ 为第 $t$ 时隙的到达数，则平均到达率为：
\[
\mathbb{E}[X_t] = \pi_{\mathrm{on}} \lambda_1 + \pi_{\mathrm{off}} \lambda_0.
\]
记 $\Delta = \lambda_1 - \lambda_0$。

\paragraph{状态的方差及自相关}
状态的方差为：
\[
\mathrm{Var}(S_t) = \pi_{\mathrm{on}} (1 - \pi_{\mathrm{on}}),
\]
状态序列的自相关系数为：
\[
\rho = 1 - \sigma_0 - \sigma_1, \quad
\mathrm{Cov}(S_t, S_{t+\tau}) = \pi_{\mathrm{on}} (1 - \pi_{\mathrm{on}}) \, \rho^{\tau}.
\]

\paragraph{到达过程的方差}
由于条件独立性（在已知状态下，不同时隙的到达独立），到达序列的方差为：
\[
\mathrm{Var}(X_t) = (\lambda_1 - \lambda_0)^2 \, \pi_{\mathrm{on}} (1 - \pi_{\mathrm{on}}) 
+ \pi_{\mathrm{on}} \lambda_1 (1 - \lambda_1) + \pi_{\mathrm{off}} \lambda_0 (1 - \lambda_0).
\]
其中第一项表示由状态切换引起的聚集（突发性），后两项为条件到达的本征波动。

\paragraph{离散度指数（ID）}
按定义，离散度指数为：
\[
\mathrm{ID} = \frac{\mathrm{Var}(X_t)}{\mathbb{E}[X_t]}.
\]
将上式代入，即得到 ON--OFF 模型下 ID 的闭式表达：
\[
\mathrm{ID} = \frac{
	(\lambda_1 - \lambda_0)^2 \, \pi_{\mathrm{on}} (1 - \pi_{\mathrm{on}}) 
	+ \pi_{\mathrm{on}} \lambda_1 (1 - \lambda_1) + \pi_{\mathrm{off}} \lambda_0 (1 - \lambda_0)
}{
	\pi_{\mathrm{on}} \lambda_1 + \pi_{\mathrm{off}} \lambda_0
}.
\]

\paragraph{特例：二值 ON/OFF 流量}
若每时隙最多到达一个包，即 $\lambda_1=1$，$\lambda_0=0$，则条件方差为零，方差简化为：
\[
\mathrm{Var}(X_t) = \pi_{\mathrm{on}} (1 - \pi_{\mathrm{on}}), \qquad
\mathbb{E}[X_t] = \pi_{\mathrm{on}}.
\]
此时
\[
\mathrm{ID} = \frac{\mathrm{Var}(X_t)}{\mathbb{E}[X_t]} = 1 - \pi_{\mathrm{on}} = \pi_{\mathrm{off}}.
\]
因此单时隙尺度下，二值 ON--OFF 流量的 ID 不会超过 1；如需体现强突发性，可通过窗口聚合（binning）或使用其他指标（如连续 ON 段长度分布、Fano 因子随窗口演化）刻画。

